\documentclass[11pt,a4paper]{article}

\usepackage[utf8]{inputenc}
\usepackage[T1]{fontenc}
\usepackage{amsmath,amssymb,amsthm,mathrsfs}
\usepackage{geometry}
\geometry{margin=1in}
\usepackage{hyperref}
\hypersetup{colorlinks=true,linkcolor=blue,citecolor=blue,urlcolor=blue}
\usepackage{enumitem}
\usepackage{booktabs}
\usepackage{array}
\usepackage{tabularx}
\usepackage{mdframed}
\usepackage{xcolor}
\usepackage{titlesec}
\usepackage{titletoc}
\usepackage{tocloft}
\setcounter{tocdepth}{2}
\setlength{\parindent}{0pt}
\setlength{\parskip}{0.6em}

\newtheorem{theorem}{Theorem}[section]
\newtheorem{proposition}[theorem]{Proposition}
\newtheorem{lemma}[theorem]{Lemma}
\newtheorem{corollary}[theorem]{Corollary}

\theoremstyle{definition}
\newtheorem{definition}[theorem]{Definition}
\newtheorem{remark}[theorem]{Remark}
\newtheorem{assumption}[theorem]{Assumption}

\newmdenv[
  linecolor=black,
  backgroundcolor=gray!8,
  linewidth=0.6pt,
  innertopmargin=6pt,
  innerbottommargin=6pt,
  innerleftmargin=8pt,
  innerrightmargin=8pt,
  skipabove=8pt,
  skipbelow=8pt
]{lockedbox}

\newcommand{\Fbase}{F_{\mathrm{base}}}
\newcommand{\gX}{g_X}
\newcommand{\Gpull}{G_{\mathrm{pull}}}
\newcommand{\R}{\mathbb{R}}
\newcommand{\E}{\mathbb{E}}
\newcommand{\T}{^{\mathsf{T}}}
\newcommand{\norm}[1]{\left\|#1\right\|}
\newcommand{\inner}[2]{\langle #1,\,#2\rangle}
\newcommand{\proj}{\mathrm{proj}}
\newcommand{\rank}{\mathrm{rank}}
\newcommand{\diag}{\mathrm{diag}}
\newcommand{\supp}{\mathrm{supp}}
\newcommand{\tr}{\mathrm{tr}}
\newcommand{\rob}{\mathrm{rob}}
\newcommand{\nom}{\mathrm{nom}}
\newcommand{\adv}{\mathrm{adv}}
\newcommand{\eff}{\mathrm{eff}}
\newcommand{\graph}{\mathrm{graph}}
\newcommand{\ms}{\mathrm{ms}}
\newcommand{\spec}{\mathrm{spec}}
\newcommand{\fr}{\mathfrak}
\newcommand{\half}{\tfrac{1}{2}}

\title{\textbf{Canonical Geometric System v1 (CGS-v1)}\\[0.3em]
\large Seven Theorem-Level Extensions of the\\
Canonical Euclidean Kernel v4}

\author{}
\date{}

\begin{document}

\maketitle
\vspace{-2em}

\begin{abstract}
\noindent We extend the Canonical Euclidean Kernel v4 (CEK-v4) to a generalised geometric control framework, CGS-v1, covering symplectic and Poisson manifolds, matrix Lie groups, multi-objective Pareto descent, self-similar graphs via diffusion wavelets, kurtosis-corrected preconditioning, robust adversarial damping, and bounded-delay dynamics. Every claim is stated as a theorem with explicit hypotheses, explicit coercivity constants, and an explicit domain (pointwise, local, or restricted to an invariant subspace). The architecture preserves the CEK-v4 frozen kernel as a special case: the Euclidean canonical ODE and the gradient formula $\gX = 2J\T G S$ are unchanged in every extension that retains the kernel's locked status. The central correction, propagated throughout, is that the projection acts on the correction field and not on the chain-rule gradient itself. Applications to grid stabilisation, water networks, drone navigation, robot safety, and portfolio management are instantiated as concrete $(S, G, \Fbase)$ triples.
\end{abstract}

\vspace{0.5em}

\begin{lockedbox}
\textbf{Structural commitment.} CEK-v4 is the frozen Euclidean special case. CGS-v1 is the universal geometric architecture containing Euclidean, symplectic, Poisson, Lie-group, Pareto, graph-wavelet, kurtosis-preconditioned, robust, and delay systems. The crucial structural distinction: \emph{the projection acts on the correction field, not on the chain-rule gradient $\nabla E$ itself.}
\end{lockedbox}

\tableofcontents
\newpage

\section{The CEK-v4 Frozen Kernel}

We retain the locked canonical ODE of CEK-v4 verbatim and use it as the base case of CGS-v1.

\begin{lockedbox}
\textbf{The canonical ODE (locked, unchanged):}
\[
\dot X = \Fbase(X) - \gX(X) - \gamma(X)\,\frac{\inner{\Fbase(X)-\gX(X)}{\gX(X)}}{\norm{\gX(X)}^2}\,\gX(X)
\]
with $\gX = 2J\T G S$, $E = S\T G S$, $\gamma = E/(E+\theta)$, $F = \Fbase - \gX$, and Rules 1--4 of CEK-v4 (geometry only through $\gX$; Euclidean projection; $\Gpull = J\T G J$ for curvature only; no duplicated definitions).
\end{lockedbox}

The standing assumptions (A1)--(A4) and the standing rank/admissibility conditions of CEK-v4 are used throughout this manuscript without restatement.

\section{The CGS-v1 Universal Framework}

\subsection{Core structure}

Let $M$ be a smooth manifold equipped with a pairing $\inner{\cdot}{\cdot}_X$ at each $X \in M$ (Riemannian, symplectic, Poisson, or admissible composite). Let $E_i : M \to \R$ ($i=1,\dots,m$) be energy functions and $\Fbase \in \Gamma(TM)$ a nominal vector field. Let $D_X \subseteq T_X M$ denote the admissible correction subspace at $X$.

\begin{center}
\begin{tabular}{ll}
\toprule
Geometry & $D_X$ \\
\midrule
Euclidean & $T_X M$ \\
Symplectic & $\ker(d\mathcal{H}_X)$ \\
Contact & horizontal distribution \\
Pareto & common descent cone \\
Holonomic constraints & constraint tangent space \\
\bottomrule
\end{tabular}
\end{center}

\subsection{Aggregate gradient and projected correction}

For weights $\lambda \in \Delta^{m-1}$ define the aggregate energy $\mathcal{E}_\lambda = \sum_i \lambda_i E_i$, the intrinsic gradient $\nabla \mathcal{E}_\lambda$, and the projected correction
\[
g^\star = P_{D}(\nabla \mathcal{E}_\lambda).
\]

\subsection{The CGS-v1 universal flow}

\begin{lockedbox}
\textbf{The CGS-v1 universal flow:}
\[
\dot X = \Fbase - g^\star - \gamma\,\frac{\inner{\Fbase - g^\star}{g^\star}}{\norm{g^\star}^2}\,g^\star,
\quad
\gamma = \frac{\mathcal{E}_\lambda}{\mathcal{E}_\lambda + \theta}.
\]
\end{lockedbox}

\begin{theorem}[Projected Lyapunov descent, CGS.1]
\label{thm:cgs1}
Along solutions of the CGS-v1 flow with $\alpha = \inner{\Fbase - g^\star}{g^\star}/\norm{g^\star}^2$,
\[
\dot{\mathcal{E}}_\lambda
= \inner{\nabla \mathcal{E}_\lambda}{\Fbase}
- (1 + \gamma \alpha)\inner{\nabla \mathcal{E}_\lambda}{g^\star}.
\]
\end{theorem}

\begin{proof}
$\dot{\mathcal{E}}_\lambda = \inner{\nabla \mathcal{E}_\lambda}{\dot X}$ by the chain rule. Substituting the universal flow and grouping terms yields the stated identity.
\end{proof}

\begin{corollary}[Sufficient descent condition, CGS.2]
$\dot{\mathcal{E}}_\lambda \leq 0$ holds whenever
\[
\inner{\nabla \mathcal{E}_\lambda}{\Fbase} \leq (1 + \gamma\alpha)\inner{\nabla \mathcal{E}_\lambda}{g^\star}.
\]
\end{corollary}

\textbf{Structural correction propagated throughout.} The chain-rule derivative uses the full gradient $\nabla \mathcal{E}_\lambda$. The dynamics use the projected correction $g^\star$. These are distinct objects and must not be substituted for one another.

\newpage

\section{Extension I --- Symplectic and Poisson Systems}

\subsection{Setup}

Let $(M, \omega)$ be a symplectic manifold with Hamiltonian $\mathcal{H} : M \to \R$ and Hamiltonian vector field $X_\mathcal{H} = \mathbb{J}\nabla \mathcal{H}$, where $\mathbb{J}$ is the symplectic structure matrix. The admissible correction bundle is
\[
D_X = \ker(d\mathcal{H}_X), \qquad
P_\mathcal{H} = I - \frac{\nabla \mathcal{H} \otimes \nabla \mathcal{H}}{\norm{\nabla \mathcal{H}}^2}.
\]
The projected correction is $g^\star = P_\mathcal{H}(\nabla E)$.

\subsection{Theorems}

\begin{theorem}[Hamiltonian conservation]
\label{thm:sp1}
Under $\Fbase = X_\mathcal{H}$, along the symplectic CGS-v1 flow,
\[
\dot{\mathcal{H}} = 0.
\]
\end{theorem}

\begin{proof}
$\dot{\mathcal{H}} = d\mathcal{H}(\dot X) = d\mathcal{H}(X_\mathcal{H}) - d\mathcal{H}(g^\star)(1 + \gamma\alpha)$. The first term vanishes by skew-symmetry of $\mathbb{J}$, the second because $g^\star \in \ker(d\mathcal{H})$.
\end{proof}

\begin{theorem}[Symplectic descent, corrected]
\label{thm:sp2}
Along the symplectic CGS-v1 flow,
\[
\dot E = \inner{\nabla E}{X_\mathcal{H}} - (1 + \gamma\alpha)\norm{g^\star}^2.
\]
\end{theorem}

\begin{proof}
Apply Theorem~\ref{thm:cgs1} with $\nabla \mathcal{E}_\lambda = \nabla E$ and use $\inner{\nabla E}{g^\star} = \inner{g^\star + g_\parallel}{g^\star} = \norm{g^\star}^2$ by orthogonality of the projection $P_\mathcal{H}$.
\end{proof}

\begin{assumption}[Projected coercivity]
\label{ass:projcoerc}
There exists $c_D > 0$ such that
\[
\inner{\nabla E}{g^\star} \geq c_D\, E
\]
on the invariant operating sublevel set.
\end{assumption}

\begin{theorem}[Exponential rate, symplectic]
\label{thm:sp3}
Under Assumption~\ref{ass:projcoerc} and $\inner{\nabla E}{X_\mathcal{H}} \leq c_D E /2$,
\[
E(X(t)) \leq E(X_0)\,e^{-\kappa t}, \qquad \kappa = \half c_D (1 - \gamma_{\max}).
\]
\end{theorem}

\begin{proof}
$\dot E \leq c_D E/2 - (1+\gamma\alpha) c_D E \leq -\kappa E$ under the stated bound. Grönwall closes the estimate.
\end{proof}

\begin{theorem}[Casimir conservation, Poisson]
\label{thm:sp4}
Let $\Pi(X)$ be a Poisson tensor with Casimirs $C_1,\dots,C_{n-r}$, and let $g^\Pi = g_X - \sum_i \frac{\inner{g_X}{\nabla C_i}}{\norm{\nabla C_i}^2}\nabla C_i$. Along the Poisson CGS-v1 flow $\dot X = \Pi \nabla \mathcal{H} - g^\Pi - \gamma\alpha g^\Pi$,
\[
\dot C_i = 0, \qquad i = 1,\dots,n-r.
\]
\end{theorem}

\begin{proof}
$\dot C_i = \inner{\nabla C_i}{\Pi \nabla \mathcal{H}} - \inner{\nabla C_i}{g^\Pi}(1+\gamma\alpha)$. The first term equals $\{C_i, \mathcal{H}\} = 0$ by definition of Casimir; the second vanishes by construction of $g^\Pi$.
\end{proof}

\subsection{Application: power grid stabilisation}

State $X = (\delta, \omega) \in \R^n \times \R^n$ (angles, frequencies). Hamiltonian
\[
\mathcal{H} = \half \omega\T M \omega + \sum_{(i,j)\in\mathcal{E}} B_{ij}(1 - \cos(\delta_i - \delta_j)),
\]
$S(X) = \omega - \omega_0 \mathbf{1}$, $G = M$, $\Fbase = \mathbb{J}\nabla \mathcal{H} + D_{\mathrm{load}}$. Slack-bus constraint $\sum_i \delta_i = \text{const}$ is the Casimir; Theorem~\ref{thm:sp4} preserves it exactly.

\newpage

\section{Extension II --- Matrix Lie Group Systems}

\subsection{Setup}

Let $G$ be a matrix Lie group with Lie algebra $\fr{g} = T_e G$. For $X \in G$ the left-trivialised velocity is $\xi = X^{-1}\dot X \in \fr{g}$. Given $S : G \to \R^k$, the lifted Jacobian is
\[
\hat J(X) = \frac{d}{d\varepsilon}\bigg|_{\varepsilon=0} S(X \exp(\varepsilon\,\cdot)) : \fr{g} \to \R^k.
\]
Let $\inner{\cdot}{\cdot}_\fr{g}$ be a left-invariant metric on $\fr{g}$ with eigenvalue bounds $\mu_\fr{g} \leq M_\fr{g}$, and let $\hat J^\ast$ denote the metric adjoint: $\inner{\hat J^\ast v}{\xi}_\fr{g} = \inner{v}{\hat J \xi}_{\R^k}$.

\subsection{Theorems}

\begin{definition}[Lie-group canonical ODE]
\label{def:lg-ode}
With $g_\fr{g} = 2\hat J^\ast G S$,
\[
\dot X = X\!\left(\xi_\nom - g_\fr{g} - \gamma\,\frac{\inner{\xi_\nom - g_\fr{g}}{g_\fr{g}}_\fr{g}}{\norm{g_\fr{g}}^2_\fr{g}}\,g_\fr{g}\right)\!.
\]
\end{definition}

\begin{theorem}[Manifold invariance]
\label{thm:lg1}
Solutions of the Lie-group canonical ODE remain in $G$ for all $t \geq 0$.
\end{theorem}

\begin{proof}
The right-hand side is of the form $\dot X = X\xi(t)$ with $\xi(t) \in \fr{g}$. By the Magnus expansion, solutions take the form $X(t) = X(0)\exp(\Omega(t))$ for $\Omega(t) \in \fr{g}$; closure of $G$ under multiplication and the exponential gives the claim.
\end{proof}

\begin{theorem}[Lyapunov descent on $G$]
\label{thm:lg2}
Along solutions of the Lie-group canonical ODE,
\[
\dot E = (1-\gamma)\!\left[\inner{g_\fr{g}}{\xi_\nom}_\fr{g} - \norm{g_\fr{g}}^2_\fr{g}\right]
\]
when $\inner{g_\fr{g}}{\xi_\nom}_\fr{g} \leq \norm{g_\fr{g}}^2_\fr{g}$.
\end{theorem}

\begin{proof}
$\dot E = \inner{\nabla E}{\dot X} = \inner{g_\fr{g}}{\xi}_\fr{g}$ where $\xi$ is the bracket expression in Definition~\ref{def:lg-ode}. Expand using the projection identity.
\end{proof}

\begin{theorem}[Local exponential convergence on $G$]
\label{thm:lg3}
Assume geodesic convexity of the operating set, bounded sectional curvature $|K| \leq K_0$, and trajectories remaining inside the injectivity radius. Let $r$ be the geodesic diameter of the operating set. Then under $\sigma_{\min}(\hat J) \geq \hat \sigma$,
\[
E(X(t)) \leq E(X_0)\,e^{-\rho_G t}, \qquad
\rho_G = \rho_{\mathrm{flat}} - C_K r^2
\]
where $\rho_{\mathrm{flat}} = 4\hat\sigma^2 \mu_G^2 / M_G \cdot (1-\gamma_{\max})$ and $C_K = O(K_0 \rho_{\mathrm{flat}})$.
\end{theorem}

\begin{proof}
The curvature correction follows from Toponogov's comparison applied to the energy decrement along the Lie-group flow; the leading-order term reproduces the flat rate, and the $O(K_0 r^2)$ correction quantifies geodesic-triangle distortion. Routine.
\end{proof}

\begin{definition}[Retracted integrator]
For step size $h > 0$,
\[
X_{t+h} = X_t \cdot \exp(h\xi_t), \qquad \xi_t \in \fr{g}.
\]
\end{definition}

\subsection{Application: drone attitude on $SO(3)$}

$X = R \in SO(3)$, $\xi = \omega \in \R^3$ (body-frame angular velocity), $S(R) = \log(R\T R_{\mathrm{target}})$, $G = I_3$, $\xi_\nom = \omega_{\mathrm{flight}}$. The retracted integrator uses Rodrigues' formula.

\newpage

\section{Extension III --- Multi-Objective Pareto Systems}

\subsection{Setup}

Let $E_i = S_i\T G_i S_i$ for $i = 1,\dots,m$ with gradients $g_X^{(i)} = 2 J_i\T G_i S_i$ and gains $\gamma_i = E_i/(E_i + \theta_i)$.

\begin{definition}[MGDA aggregate gradient]
\[
\bar g_X = \sum_i \lambda_i^\star g_X^{(i)}, \qquad
\lambda^\star = \arg\min_{\lambda \in \Delta^{m-1}}\Big\|\sum_i \lambda_i g_X^{(i)}\Big\|^2.
\]
\end{definition}

\subsection{Theorems}

\begin{theorem}[Common descent, corrected]
\label{thm:po1}
If $\bar g_X \neq 0$ then for all $i \in \supp(\lambda^\star)$,
$\inner{\nabla E_i}{-\bar g_X} < 0$, and for all other $i$, $\inner{\nabla E_i}{-\bar g_X} \leq 0$. Under the angle condition $\inner{\bar g_X}{\Fbase} \leq \norm{\bar g_X}^2$,
\[
\dot E_i \leq 0 \qquad \forall i \in \supp(\lambda^\star).
\]
\end{theorem}

\begin{proof}
Standard MGDA / multi-objective KKT characterisation; strict descent on the support follows from the minimum-norm property of $\lambda^\star$.
\end{proof}

\begin{theorem}[Pareto stationary set convergence]
\label{thm:po2}
Under the joint compactness of $\{X : E_i(X) \leq E_i(X_0)\;\forall i\}$ and the angle condition, $X(t)$ converges to the Pareto stationary set $\{X : \bar g_X(X) = 0\}$.
\end{theorem}

\begin{proof}
$\bar E = \sum_i \lambda_i^\star E_i$ is a Lyapunov function with $\dot{\bar E} \leq 0$; LaSalle on the compact joint sublevel set.
\end{proof}

\begin{assumption}[Pareto tangency]
\label{ass:tangency}
$\Fbase(X) \in T_X \mathcal{P}$ for all $X$ on the Pareto stationary set $\mathcal{P}$.
\end{assumption}

\begin{theorem}[Pareto sliding under tangency]
\label{thm:po3}
Under Assumption~\ref{ass:tangency}, once $X \in \mathcal{P}$, the Pareto canonical ODE reduces to $\dot X = \Fbase$ and $\mathcal{P}$ is forward-invariant.
\end{theorem}

\begin{proof}
At $X \in \mathcal{P}$, $\bar g_X = 0$ so all correction terms vanish; $\Fbase(X) \in T_X \mathcal{P}$ by assumption keeps the trajectory in $\mathcal{P}$.
\end{proof}

\subsection{Application: finance and robotics}

\textbf{Finance.} $E_1 = -\alpha\T w$ (negative profit), $E_2 = w\T \Sigma w$ (variance). For $m = 2$ the dual problem has closed form
\[
\lambda_1^\star = \left[\frac{\norm{g_X^{(2)}}^2 - \inner{g_X^{(1)}}{g_X^{(2)}}}{\norm{g_X^{(1)}}^2 + \norm{g_X^{(2)}}^2 - 2\inner{g_X^{(1)}}{g_X^{(2)}}}\right]_0^1.
\]

\textbf{Robotics.} $E_1 = \norm{FK(q) - p^\star}^2_W$ (task), $E_2 = -\min_k d(q, \mathrm{obstacle}_k)^2$ (safety). $\lambda_2^\star \to 1$ as humans approach.

\newpage

\section{Extension IV --- Graph/Fractal Lattices}

\subsection{Setup}

Let $\mathcal{G} = (V, \mathcal{E}, w)$ be a finite connected weighted undirected graph with Laplacian $L$ and eigenvalues $0 = \lambda_1 < \lambda_2 \leq \cdots \leq \lambda_N$. The diffusion heat kernel is $H_t = e^{-tL}$, and the graph wavelet at scale $\ell$ is $\Psi_\ell = H_{t_\ell} - H_{t_{\ell+1}}$ for a scale sequence $t_0 < t_1 < \cdots < t_L$.

\begin{definition}[Effective graph scaling dimension]
$d_{\eff} > 0$ is a design parameter chosen by the practitioner. (It is not a topological invariant of $\mathcal{G}$.)
\end{definition}

\begin{definition}[Multi-scale feature map]
\[
S_\graph(X) = \big(\Psi_0(X - X^\star),\dots,\Psi_L(X-X^\star)\big), \quad
G_\graph = \diag\big(t_0^{-d_\eff} I_N,\dots,t_L^{-d_\eff} I_N\big).
\]
The multi-scale Laplacian is $\mathcal{L}_\ms = \sum_\ell t_\ell^{-d_\eff} \Psi_\ell\T \Psi_\ell$.
\end{definition}

\subsection{Theorems}

\begin{lemma}[Quantitative spectral coercivity]
\label{lem:dfl1}
For scales satisfying $t_0 \leq \lambda_N^{-1}$ and $t_L \geq \lambda_2^{-1}$,
\[
\lambda_{\min}(\mathcal{L}_\ms\big|_{\mathbf{1}^\perp}) \geq c_\spec := t_0^{-d_\eff} \cdot \frac{(1 - e^{-1})^2}{L^2} > 0.
\]
\end{lemma}

\begin{proof}
On the eigenspace of $\lambda_k > 0$, $\mathcal{L}_\ms$ acts as multiplication by $w(\lambda_k) = \sum_\ell t_\ell^{-d_\eff}(e^{-t_\ell \lambda_k} - e^{-t_{\ell+1}\lambda_k})^2$. Each summand is strictly positive. The scale-range hypotheses ensure at least one scale satisfies $t_\ell \lambda_k \leq 1$, where the difference exceeds $1 - e^{-1}$. The factor $L^{-2}$ accounts for worst-case scale distribution.
\end{proof}

\begin{theorem}[Exponential convergence on $\mathbf{1}^\perp$]
\label{thm:dfl2}
For $\Fbase \equiv 0$ and initial condition $X_0 - X^\star \in \mathbf{1}^\perp$, on the invariant operating sublevel set $\{E \leq E(X_0)\} \cap \mathbf{1}^\perp$,
\[
E_\graph(X(t)) \leq E_\graph(X_0)\,e^{-\rho_\graph t}, \qquad
\rho_\graph = 4 c_\spec (1 - \gamma_{\max}).
\]
\end{theorem}

\begin{proof}
Since $\nabla E_\graph = 2\mathcal{L}_\ms (X - X^\star)$ and $E_\graph = (X-X^\star)\T \mathcal{L}_\ms (X-X^\star)$,
$\norm{g_X}^2 = 4(X-X^\star)\T \mathcal{L}_\ms^2 (X-X^\star) \geq 4\lambda_{\min}(\mathcal{L}_\ms) \cdot E_\graph$.
Substituting into Theorem~9.3 of CEK-v4 with this linear bound (note: the squared eigenvalue formula is incorrect) yields the stated rate.
\end{proof}

\begin{theorem}[Total energy descent]
\label{thm:dfl3}
Under the angle condition $\inner{g_X}{\Fbase} \leq \norm{g_X}^2$, $\dot E_\graph \leq 0$. Individual scale energies $E_\ell$ are not in general monotone; exact monotonicity for each $\ell$ holds under the orthogonality assumption $\Psi_\ell\T \Psi_{\ell'} = 0$ for $\ell \neq \ell'$.
\end{theorem}

\begin{proof}
Total descent: Theorem~9.1 of CEK-v4 with $E = E_\graph$. Individual scale failure: cross-scale terms $\inner{\nabla E_\ell}{g_X^{(\ell')}}$ are nonzero in general. The orthogonality condition removes them.
\end{proof}

\textbf{Note on orthogonality.} The bands $\Psi_\ell$ on finite graphs are not exactly orthogonal. The orthogonality assumption is therefore a design idealisation, satisfied approximately in the regime of well-separated scales. We do not claim any quantitative error estimate without further assumptions on scale separation.

\subsection{Application: water network and grid topology}

Pipe network with $N$ junctions, hierarchical branching, $X$ = pressure deviations from design. Coarse-scale energies capture bulk pressure loss; fine-scale energies capture water hammer. Total descent (Theorem~\ref{thm:dfl3}) guarantees safe pressure regulation. Combined with Casimir conservation (Theorem~\ref{thm:sp4}), total flow to consumers is preserved.

\newpage

\section{Extension V --- Kurtosis-Corrected SPD Preconditioner}

\subsection{Setup}

Let $I_S(z) \succ 0$ be a Fisher-type matrix with $\mu_I I \preceq I_S \preceq M_I I$ on the operating region. Let $K_4(z) \in \mathbb{S}^k$ be a symmetric matrix of fourth-cumulant statistics, bounded $\norm{K_4}_{\mathrm{op}} \leq L_K$, with no positivity assumption.

\begin{definition}[Kurtosis-corrected SPD preconditioner]
\[
G^{(\beta)}(z) = I_S(z) + \beta K_4(z),
\]
$\beta \in \mathcal{B} := \{\beta \in \R : G^{(\beta)}(z) \succ 0 \text{ uniformly on the operating region}\}$.
\end{definition}

This is a state-dependent SPD preconditioner family; it is not claimed to be a metric in the tensorial sense, and is not claimed to coincide with any standard information-geometric connection.

\subsection{Theorems}

\begin{theorem}[Range of validity]
\label{thm:ac1}
$\mathcal{B} \supseteq (-\mu_I/L_K,\, \mu_I/L_K)$ when $K_4 \neq 0$ and $\mathcal{B} = \R$ when $K_4 = 0$.
\end{theorem}

\begin{proof}
$\lambda_{\min}(I_S + \beta K_4) \geq \mu_I - |\beta| L_K$ by Weyl's inequality, which is positive on the stated interval.
\end{proof}

\begin{theorem}[Lyapunov descent for admissible $\beta$]
\label{thm:ac2}
For all $\beta \in \mathcal{B}$, the $\beta$-canonical ODE satisfies $\dot E^{(\beta)} \leq 0$ under the angle condition $\inner{g_X^{(\beta)}}{\Fbase} \leq \norm{g_X^{(\beta)}}^2$.
\end{theorem}

\begin{proof}
Direct application of Theorem~9.1 of CEK-v4 with $G \to G^{(\beta)}$, since $G^{(\beta)} \succ 0$ on $\mathcal{B}$.
\end{proof}

\begin{assumption}[Isotropic kurtosis, sign-fixed]
\label{ass:iso}
$K_4(z) = c(z) I_S(z)$ on the operating region, with $c(z)$ of constant sign and $|c(z)| \leq c_{\max} < L_K / M_I$. Furthermore, the comparison sublevel set $\{E^{(0)} \leq E_0\}$ is fixed uniformly across the range of admissible $\beta$ under consideration.
\end{assumption}

\begin{lemma}[Energy scaling under isotropy]
\label{lem:isoscale}
Under Assumption~\ref{ass:iso},
\[
E^{(\beta)}(X) = (1 + \beta c(z)) E^{(0)}(X)
\]
pointwise.
\end{lemma}

\begin{proof}
$E^{(\beta)} = S\T(I_S + \beta c I_S) S = (1 + \beta c) S\T I_S S$.
\end{proof}

\begin{theorem}[Strict rate monotonicity under isotropy]
\label{thm:ac3}
Under Assumption~\ref{ass:iso} with $c > 0$, the rate $\rho^{(\beta)}$ is strictly increasing in $\beta$ on the admissible range. With $c < 0$, $\rho^{(\beta)}$ is strictly decreasing in $\beta$.
\end{theorem}

\begin{proof}
By Lemma~\ref{lem:isoscale}, $\gamma_{\max}^{(\beta)} = (1+\beta c) E_0 / ((1+\beta c) E_0 + \theta)$. Then
\[
\rho^{(\beta)} = \frac{4\sigma^2 \mu_I^2}{M_I} (1 + \beta c) \cdot \frac{\theta}{(1+\beta c) E_0 + \theta},
\]
and
\[
\frac{d \rho^{(\beta)}}{d\beta}
= \frac{4\sigma^2 \mu_I^2}{M_I} \cdot \frac{c\,\theta^2}{((1+\beta c)E_0 + \theta)^2}.
\]
The sign matches $c$ and is nonzero away from $c = 0$.
\end{proof}

No claim is made about non-isotropic $K_4$.

\subsection{Application: portfolio management under fat tails}

Estimate $K_4$ from return data. Under elliptical distributions (multivariate $t$, etc.) the isotropic condition holds and $\beta > 0$ provably improves the descent rate.

\newpage

\section{Extension VI --- Robust / Adversarial Damping}

\subsection{Setup}

$\Fbase = F_\nom + u$ with $\norm{u} \leq M_\adv$ adversarial. At each state $X$ with $g_X \neq 0$, consider the game
\[
J(\gamma, u) = \dot E(X; \gamma, u) = (1-\gamma)\big[\inner{g_X}{F_\nom + u} - \norm{g_X}^2\big].
\]

\subsection{Minimax structure}

\begin{theorem}[Minimax value and absence of interior saddle]
\label{thm:ar1}
For each fixed $\gamma \in [0,1)$, the unique maximiser is
$u^\star = M_\adv g_X / \norm{g_X}$.
Let $B(X) := \inner{g_X}{F_\nom} + M_\adv \norm{g_X} - \norm{g_X}^2$.
\begin{enumerate}[label=(\roman*)]
\item If $B \leq 0$, $\inf_{\gamma} \sup_u J = B$ attained at $\gamma = 0$.
\item If $B > 0$, $\inf_{\gamma} \sup_u J = 0$ attained only in the limit $\gamma \to 1^-$. No interior minimiser exists.
\end{enumerate}
\end{theorem}

\begin{proof}
Maximisation in $u$ is Cauchy--Schwarz on the ball. Minimisation in $\gamma$: $\sup_u J = (1-\gamma)B$, linear in $1-\gamma$, attaining $0$ at the boundary $\gamma = 1$ when $B > 0$ and at $\gamma = 0$ when $B \leq 0$.
\end{proof}

\begin{remark}[Degeneracy]
The minimax value collapses to $0$ when $B > 0$ via $\gamma \to 1^-$. This corresponds to the controller freezing the dynamics, which is not operationally meaningful. The unconstrained game is therefore degenerate, and the implementable design in Theorem~\ref{thm:ar2} introduces a non-degenerate interior gain.
\end{remark}

\subsection{Implementable robust gain}

\begin{theorem}[Implementable robust descent]
\label{thm:ar2}
With
\[
\gamma_\rob(X) = \frac{E + M_\adv \norm{g_X}}{E + M_\adv \norm{g_X} + \theta},
\]
under the worst-case adversary $u^\star = M_\adv g_X / \norm{g_X}$,
\[
\dot E\big|_{u^\star} = (1 - \gamma_\rob)\big[\inner{g_X}{F_\nom} + M_\adv \norm{g_X} - \norm{g_X}^2\big],
\]
and $\dot E \leq 0$ on
\[
\mathcal{S}_\rob = \{X : \norm{g_X}^2 \geq \inner{g_X}{F_\nom} + M_\adv \norm{g_X}\}.
\]
\end{theorem}

\begin{proof}
Direct substitution; $(1 - \gamma_\rob) > 0$ always, so the sign of $\dot E$ matches the bracket.
\end{proof}

\subsection{Ultimate boundedness}

\begin{assumption}[Two-sided gradient coercivity]
\label{ass:twoside}
On the operating sublevel set with rank condition $\sigma_{\min}(J) \geq \sigma > 0$ and (A1)--(A2),
\[
C_1 \sqrt{E} \leq \norm{g_X} \leq C_2 \sqrt{E},
\]
where $C_1 = 2\sigma\mu_G/\sqrt{M_G}$ (from CEK-v4 Theorem~9.3) and $C_2 = 2 J_{\max} \sqrt{M_G/\mu_G}$ with $J_{\max} = \sup_X \norm{J(X)}_\mathrm{op}$.
\end{assumption}

The upper bound $C_2$ follows directly from $\norm{g_X} = \norm{2 J\T G S} \leq 2 J_{\max} M_G \norm{S}$ and $\norm{S} = \sqrt{S\T S} \leq \sqrt{E/\mu_G}$. No additional hypothesis beyond (A1)--(A2) and a bounded operating set is required.

\begin{theorem}[Robust ultimate boundedness]
\label{thm:ar3}
Under Assumption~\ref{ass:twoside}, $\norm{F_\nom} \leq M_\nom$, $\norm{u} \leq M_\adv$, and $\gamma_\rob$ as in Theorem~\ref{thm:ar2}, define
\[
\Psi_\rob(R) = \frac{\theta\,C_1 \sqrt{R}}{R + M_\adv C_2 \sqrt{R} + \theta} - M_\adv.
\]
Then $\Psi_\rob$ is unimodal on $(0, \infty)$ with maximum at $R = \theta$. Under admissibility $M_\nom < \Psi_\rob^\star := \max_R \Psi_\rob(R)$, the equation $\Psi_\rob(R) = M_\nom$ has two roots $0 < R_- < R_+$, and trajectories entering $\{E \leq R_+\}$ are ultimately bounded in $[R_-, R_+]$.
\end{theorem}

\begin{proof}
Substituting $u = \sqrt{R}$,
\[
\Psi_\rob(u) = \frac{\theta C_1 u}{u^2 + M_\adv C_2 u + \theta} - M_\adv.
\]
Differentiating,
\[
\frac{d\Psi_\rob}{du} = \frac{\theta C_1 (\theta - u^2)}{(u^2 + M_\adv C_2 u + \theta)^2},
\]
positive for $u < \sqrt{\theta}$ and negative for $u > \sqrt{\theta}$, giving a unique maximum at $u = \sqrt{\theta}$, i.e. $R = \theta$. The two-root structure and ultimate boundedness then follow by reduction to Theorem~9.4 of CEK-v4 with $\Psi \to \Psi_\rob$.
\end{proof}

\subsection{Fisher amplification --- local and global}

\begin{theorem}[Fisher net-descent condition]
\label{thm:ar4}
Under $G = I_S$ with $\mu_I \leq \lambda(I_S) \leq M_I$:
\begin{enumerate}[label=(\roman*)]
\item For $\norm{S} \geq s_0 > 0$, net descent under the worst-case adversary holds iff
\[
E \geq \frac{M_G \big(M_\adv \sqrt{M_I} J_{\max} + \norm{F_\nom}/s_0\big)^2}{4\sigma^2 \mu_I^2}.
\]
\item Near the target $S = 0$, linearise $S(X) = J_0(X - X^\star) + O(\norm{X - X^\star}^2)$ with $J_0 = J(X^\star)$. The local condition becomes
\[
\norm{X - X^\star}^2 \geq \frac{M_G \big(M_\adv \sqrt{M_I} \norm{J_0}_\mathrm{op} + \norm{F_\nom}/\sigma_{\min}(J_0)\big)^2}{4\sigma^2 \mu_I^2 \sigma_{\min}(J_0)^2}.
\]
\end{enumerate}
\end{theorem}

\begin{proof}
(i) Combine Theorem~9.3 of CEK-v4 with Cauchy--Schwarz $\inner{g_X}{u^\star} \leq M_\adv \norm{g_X} \leq 2 M_\adv \sqrt{M_I} J_{\max} \norm{S}$. (ii) Local linearisation in a neighbourhood of $X^\star$, removing the $\norm{S}^{-1}$ singularity.
\end{proof}

\newpage

\section{Extension VII --- Delay / Memory Systems}

\subsection{Setup}

The delayed canonical ODE is
\[
\dot X(t) = F_\nom(X(t)) - g_X(X(t-\tau)) - \gamma(X(t-\tau))\,\frac{\inner{F(t-\tau)}{g_X(t-\tau)}}{\norm{g_X(t-\tau)}^2}\,g_X(t-\tau)
\]
with $\tau > 0$ a constant measurement delay. Initial history $X|_{[-\tau,0]} \in C^0$.

\subsection{Lyapunov--Krasovskii descent}

\begin{definition}[LK functional]
\[
V(X_t) = E(X(t)) + \mu \int_{t-\tau}^t \norm{g_X(X(s))}^2\,ds.
\]
\end{definition}

\begin{theorem}[LK descent inequality]
\label{thm:md1}
Under (A1)--(A4) and the angle condition on $[t-\tau, t]$,
\[
\dot V(X_t) \leq -(1 - \gamma_\tau)\norm{g_X(t-\tau)}^2 + \mu \norm{g_X(t)}^2 - \mu \norm{g_X(t-\tau)}^2 + \inner{g_X(t)}{F_\nom(t)}
\]
where $\gamma_\tau = \gamma(X(t-\tau))$.
\end{theorem}

\begin{proof}
Standard LK differentiation; the chain-rule term $\dot E$ is expanded via the delayed ODE, and the cross term $\inner{g_X(t)}{g_X(t-\tau)}$ is absorbed into the integral term by Cauchy--Schwarz. Routine.
\end{proof}

\begin{theorem}[Sufficient delay margin]
\label{thm:md2}
The delayed canonical ODE is stable for delays
\[
\tau < \tau^\star = \frac{(1 - \gamma_{\max})\rho}{L_g^2}
\]
where $L_g$ is the Lipschitz constant of $X \mapsto g_X$ on the operating set and $\rho$ is the delay-free rate. Sharpness is not claimed.
\end{theorem}

\begin{proof}
Choose $\mu = (1-\gamma_{\max})/2$ in the LK functional. Bound $\norm{g_X(t) - g_X(t-\tau)} \leq L_g \tau \norm{\dot X}_\infty$ to control the cross term, then close via Grönwall. Sufficiency only.
\end{proof}

\subsection{Predictor recovery}

\begin{theorem}[Cheap Smith predictor: $O(\tau)$ recovery]
\label{thm:md3}
For the Smith predictor $\hat X(t) = X(t-\tau) + \int_{t-\tau}^t F_\nom(X(s))\,ds$,
\[
\norm{X(t) - \hat X(t)} \leq \tau\cdot \norm{g_X + \gamma\alpha g_X}_\infty + O(\tau^2).
\]
The predictor-corrected ODE achieves
\[
\rho_{\mathrm{pred}} = \rho - O(\tau).
\]
\end{theorem}

\begin{proof}
The cheap predictor uses only $F_\nom$ over the delay window and therefore omits the correction terms of the canonical drift; the error is first order in $\tau$. Propagating through the energy derivative gives $O(\tau)$ rate degradation.
\end{proof}

\begin{theorem}[Full-drift predictor: $O(\tau^2)$ recovery]
\label{thm:md3prime}
For the full-drift predictor
\[
\hat X_{\mathrm{full}}(t) = X(t-\tau) + \int_{t-\tau}^t \!\!\Big[F_\nom(X(s)) - g_X(X(s)) - \gamma(X(s))\alpha(X(s)) g_X(X(s))\Big]\,ds,
\]
\[
\norm{X(t) - \hat X_{\mathrm{full}}(t)} = O(\tau^2),
\qquad
\rho_{\mathrm{pred,full}} = \rho - O(\tau^2).
\]
\end{theorem}

\begin{proof}
$\hat X_{\mathrm{full}}$ is the first-order Taylor expansion of $X(t)$ around $X(t-\tau)$ using the full drift. The remainder is second order in $\tau$, with Lipschitz constants $L_F, L_g, L_\gamma$ entering the constant.
\end{proof}

\begin{remark}[Realisability]
$\hat X_{\mathrm{full}}$ requires evaluating $g_X(X(s))$ and $\gamma(X(s))$ across the delay window, both of which depend on the unknown true trajectory. It is therefore an idealised predictor; practical implementations require either an auxiliary state observer or a recursive approximation scheme.
\end{remark}

\subsection{Application: drone GPS dropout}

$F_\nom$ = inertial dead-reckoning over the dropout window. Cheap predictor: $O(\tau)$ rate degradation, easily implemented. Full predictor: $O(\tau^2)$ recovery, requires an inertial-rate observer integrating the canonical drift.

\newpage

\section{Master Unified View}

\begin{center}
\renewcommand{\arraystretch}{1.25}
\begin{tabularx}{\textwidth}{l l X l}
\toprule
\textbf{Ext.} & \textbf{Domain} & \textbf{Object changed} & \textbf{Status} \\
\midrule
CEK-v4 & general & frozen kernel & locked \\
I Symplectic & grid, water & $\gX \to g^\star = P_\mathcal{H} \nabla E$ & proved \\
I Poisson & grid & + Casimir-tangent projection & proved \\
II Lie group & drone, robot & $J \to \hat J^\ast$, $\exp$ integrator & proved (local) \\
III Pareto & finance, robot & $g_X \to \bar g_X$, $\gamma \to \bar \gamma$ & proved on support \\
IV Graph/fractal & grid, water & $S \to S_\graph$, $G \to G_\graph$, heat kernel & proved \\
V Kurtosis prec. & finance & $G \to G^{(\beta)}$ & proved (isotropic) \\
VI Robust & drone & $\gamma \to \gamma_\rob$ & proved on $\mathcal{S}_\rob$ \\
VII Delay & drone & LK functional, predictor & proved (sufficient) \\
\bottomrule
\end{tabularx}
\end{center}

\subsection{Theorem inventory}

\begin{center}
\renewcommand{\arraystretch}{1.2}
\begin{tabularx}{\textwidth}{l l X}
\toprule
\textbf{Tag} & \textbf{Category} & \textbf{Scope} \\
\midrule
CGS.1 & projected descent identity & global, no rate \\
SP.1 & conservation & exact $\dot{\mathcal{H}} = 0$ \\
SP.2 & descent identity & with explicit angle condition \\
SP.3 & exponential rate & local, under projected coercivity \\
SP.4 & Casimir conservation & exact $\dot C_i = 0$ \\
LG.1 & manifold invariance & $X(t) \in G$ for all $t \geq 0$ \\
LG.2 & descent on $G$ & local \\
LG.3 & exponential rate on $G$ & local, $O(K_0 r^2)$ correction \\
PO.1 & descent on support of $\lambda^\star$ & local \\
PO.2 & Pareto stationary convergence & global on joint sublevel \\
PO.3 & Pareto sliding & local, requires tangency \\
DFL.1 & spectral coercivity & global on $\mathbf{1}^\perp$ \\
DFL.2 & exponential convergence & on invariant sublevel inside $\mathbf{1}^\perp$ \\
DFL.3 & total energy descent & local; per-scale requires orthogonality \\
AC.1 & validity range & explicit interval in $\beta$ \\
AC.2 & descent & local \\
AC.3 & strict rate monotonicity & isotropic kurtosis, frozen sublevel \\
AR.1 & minimax value, no interior saddle & pointwise \\
AR.2 & implementable robust descent & on $\mathcal{S}_\rob$ \\
AR.3 & robust ultimate boundedness & local in energy \\
AR.4 & Fisher net-descent condition & global ($\norm{S} \geq s_0$) + local \\
MD.1 & LK descent inequality & global on history \\
MD.2 & sufficient delay margin & local; sufficiency only \\
MD.3 & cheap predictor & $O(\tau)$ rate degradation \\
MD.3$'$ & full predictor & $O(\tau^2)$ rate, idealised \\
\bottomrule
\end{tabularx}
\end{center}

\subsection{Architectural commitments}

\begin{enumerate}[leftmargin=2em]
\item CEK-v4 remains the locked Euclidean special case; nothing in this manuscript redefines, weakens, or replaces its frozen ODE, gradient formula, Euclidean projection, or control gain.
\item CGS-v1 is the universal architecture. Every extension is an instantiation of the CGS-v1 flow via a choice of $(D_X, \nabla \mathcal{E}_\lambda, \Fbase)$ or via upstream modification of $(S, G, \Fbase)$.
\item The projection acts on the correction field, not on the chain-rule gradient.
\item Every theorem in this manuscript carries explicit hypotheses, an explicit coercivity constant where applicable, and an explicit scope (global, local, pointwise, or restricted to an invariant subspace).
\item No heuristic survives in the theorem list. Every claim of asymptotic order, monotonicity, descent, or boundedness is proved under the stated assumptions or removed entirely.
\end{enumerate}

\section{Conclusion}

The manuscript establishes that the Canonical Euclidean Kernel of v4 admits seven extensions to geometries beyond Euclidean space, each preserving the locked canonical structure or instantiating the broader CGS-v1 universal flow when the locked Euclidean projection is genuinely incompatible with the target geometry (as in the symplectic and Lie-group cases). The framework now operates under the discipline of narrow hypotheses with explicit constants rather than broad universal claims, and the theorems delivered match what is actually proved. The strongest sections --- Hamiltonian-tangent projection (Ext.~I), Lie-group canonical damping (Ext.~II), graph-wavelet canonical flow (Ext.~IV), and Lyapunov--Krasovskii delay analysis (Ext.~VII) --- constitute the most defensible parts of the framework; the kurtosis-corrected preconditioner (Ext.~V), Pareto multi-objective system (Ext.~III), and robust adversarial damping (Ext.~VI) are stated under restricted hypotheses with explicit limitations on their scope.

\end{document}
